% Options for packages loaded elsewhere
\PassOptionsToPackage{unicode}{hyperref}
\PassOptionsToPackage{hyphens}{url}
\documentclass[
]{article}
\usepackage{xcolor}
\usepackage[margin=1in]{geometry}
\usepackage{amsmath,amssymb}
\setcounter{secnumdepth}{-\maxdimen} % remove section numbering
\usepackage{iftex}
\ifPDFTeX
  \usepackage[T1]{fontenc}
  \usepackage[utf8]{inputenc}
  \usepackage{textcomp} % provide euro and other symbols
\else % if luatex or xetex
  \usepackage{unicode-math} % this also loads fontspec
  \defaultfontfeatures{Scale=MatchLowercase}
  \defaultfontfeatures[\rmfamily]{Ligatures=TeX,Scale=1}
\fi
\usepackage{lmodern}
\ifPDFTeX\else
  % xetex/luatex font selection
\fi
% Use upquote if available, for straight quotes in verbatim environments
\IfFileExists{upquote.sty}{\usepackage{upquote}}{}
\IfFileExists{microtype.sty}{% use microtype if available
  \usepackage[]{microtype}
  \UseMicrotypeSet[protrusion]{basicmath} % disable protrusion for tt fonts
}{}
\makeatletter
\@ifundefined{KOMAClassName}{% if non-KOMA class
  \IfFileExists{parskip.sty}{%
    \usepackage{parskip}
  }{% else
    \setlength{\parindent}{0pt}
    \setlength{\parskip}{6pt plus 2pt minus 1pt}}
}{% if KOMA class
  \KOMAoptions{parskip=half}}
\makeatother
\usepackage{longtable,booktabs,array}
\usepackage{calc} % for calculating minipage widths
% Correct order of tables after \paragraph or \subparagraph
\usepackage{etoolbox}
\makeatletter
\patchcmd\longtable{\par}{\if@noskipsec\mbox{}\fi\par}{}{}
\makeatother
% Allow footnotes in longtable head/foot
\IfFileExists{footnotehyper.sty}{\usepackage{footnotehyper}}{\usepackage{footnote}}
\makesavenoteenv{longtable}
\usepackage{graphicx}
\makeatletter
\newsavebox\pandoc@box
\newcommand*\pandocbounded[1]{% scales image to fit in text height/width
  \sbox\pandoc@box{#1}%
  \Gscale@div\@tempa{\textheight}{\dimexpr\ht\pandoc@box+\dp\pandoc@box\relax}%
  \Gscale@div\@tempb{\linewidth}{\wd\pandoc@box}%
  \ifdim\@tempb\p@<\@tempa\p@\let\@tempa\@tempb\fi% select the smaller of both
  \ifdim\@tempa\p@<\p@\scalebox{\@tempa}{\usebox\pandoc@box}%
  \else\usebox{\pandoc@box}%
  \fi%
}
% Set default figure placement to htbp
\def\fps@figure{htbp}
\makeatother
\setlength{\emergencystretch}{3em} % prevent overfull lines
\providecommand{\tightlist}{%
  \setlength{\itemsep}{0pt}\setlength{\parskip}{0pt}}
\usepackage{bookmark}
\IfFileExists{xurl.sty}{\usepackage{xurl}}{} % add URL line breaks if available
\urlstyle{same}
\hypersetup{
  pdftitle={Taller Aplicado: Análisis y Predicción de Churn Bancario},
  pdfauthor={Machine Learning para Economía y Finanzas},
  hidelinks,
  pdfcreator={LaTeX via pandoc}}

\title{Taller Aplicado: Análisis y Predicción de Churn Bancario}
\author{Machine Learning para Economía y Finanzas}
\date{}

\begin{document}
\maketitle

\section{Caso de Estudio: Predicción de Churn en una Entidad
Financiera}\label{caso-de-estudio-predicciuxf3n-de-churn-en-una-entidad-financiera}

Una entidad financiera ha detectado un incremento en la tasa de abandono
de clientes durante los últimos periodos. El área de analítica ha
desarrollado un modelo de regresión logística que permite estimar la
probabilidad de que un cliente abandone el banco (churn), utilizando
información demográfica, financiera y comportamental.

La variable objetivo del estudio es:

\begin{itemize}
\tightlist
\item
  Exited = 1: El cliente abandonó el banco\\
\item
  Exited = 0: El cliente permanece activo
\end{itemize}

El propósito de este taller es analizar la estructura de los datos,
estudiar sus distribuciones, identificar relaciones bivariadas con la
variable respuesta y comprender la interpretación del modelo mediante
Odds Ratios.

\section{Contexto de las Variables del
Modelo}\label{contexto-de-las-variables-del-modelo}

Las siguientes variables forman parte del modelo predictivo:

\begin{longtable}[]{@{}
  >{\raggedright\arraybackslash}p{(\linewidth - 4\tabcolsep) * \real{0.3667}}
  >{\raggedright\arraybackslash}p{(\linewidth - 4\tabcolsep) * \real{0.4333}}
  >{\raggedright\arraybackslash}p{(\linewidth - 4\tabcolsep) * \real{0.2000}}@{}}
\toprule\noalign{}
\begin{minipage}[b]{\linewidth}\raggedright
Variable
\end{minipage} & \begin{minipage}[b]{\linewidth}\raggedright
Descripción
\end{minipage} & \begin{minipage}[b]{\linewidth}\raggedright
Tipo
\end{minipage} \\
\midrule\noalign{}
\endhead
\bottomrule\noalign{}
\endlastfoot
CreditScore & Puntaje crediticio del cliente. Indica el nivel de riesgo
financiero histórico del cliente. & Numérica \\
Geography & País de residencia del cliente (France, Germany, Spain). Se
codifica mediante variables dummy tomando France como categoría base. &
Categórica \\
Gender & Género del cliente (Female, Male). Se codifica como variable
dummy. & Categórica \\
Age & Edad del cliente en años. Variable demográfica relevante en
decisiones financieras. & Numérica \\
Tenure & Número de años que el cliente ha permanecido en el banco. &
Numérica \\
Balance & Saldo disponible en la cuenta bancaria del cliente. &
Numérica \\
NumOfProducts & Número de productos financieros que el cliente posee en
el banco. & Numérica \\
HasCrCard & Indica si el cliente posee tarjeta de crédito (0 = No, 1 =
Sí). & Binaria \\
IsActiveMember & Indica si el cliente es activo en el banco (0 = No, 1 =
Sí). & Binaria \\
EstimatedSalary & Estimación del salario anual del cliente. &
Numérica \\
\end{longtable}

\section{1. Estructura de los Datos}\label{estructura-de-los-datos}

\begin{enumerate}
\def\labelenumi{\arabic{enumi}.}
\tightlist
\item
  ¿Cuántas variables tiene la base de datos?
\item
  ¿Cuántas son cuantitativas?
\item
  ¿Cuántas son categóricas o binarias?
\item
  ¿Qué implicaciones tiene esta combinación de variables para el
  análisis estadístico?
\end{enumerate}

\section{2. Análisis Descriptivo
Univariado}\label{anuxe1lisis-descriptivo-univariado}

\begin{enumerate}
\def\labelenumi{\arabic{enumi}.}
\setcounter{enumi}{4}
\tightlist
\item
  Muestre las distribuciones de las variables que considere relevantes e
  interprete estos datos.
\end{enumerate}

\section{3. Análisis Bivariado con la Variable
Respuesta}\label{anuxe1lisis-bivariado-con-la-variable-respuesta}

\begin{enumerate}
\def\labelenumi{\arabic{enumi}.}
\setcounter{enumi}{5}
\tightlist
\item
  Compare la distribución de Age entre clientes que abandonan y los que
  permanecen.

  \begin{itemize}
  \tightlist
  \item
    ¿Se observan diferencias relevantes?
  \end{itemize}
\item
  Analice la relación entre IsActiveMember y Exited.

  \begin{itemize}
  \tightlist
  \item
    ¿Los clientes activos presentan menor abandono?
  \end{itemize}
\item
  Evalúe la relación entre NumOfProducts y Exited.

  \begin{itemize}
  \tightlist
  \item
    ¿Mayor número de productos se asocia con menor churn?
  \end{itemize}
\item
  Identifique qué variable parece tener mayor relación con el abandono,
  justificando con base en distribuciones comparativas.
\end{enumerate}

\section{4. Interpretación del Modelo: Odds
Ratios}\label{interpretaciuxf3n-del-modelo-odds-ratios}

\begin{enumerate}
\def\labelenumi{\arabic{enumi}.}
\setcounter{enumi}{9}
\tightlist
\item
  Identifique los factores de riesgo y los factores protectores del
  modelo. Recuerde: los odds ratios mayores que 1 son factores de
  riesgo; los odds ratios menores que 1 son factores protectores.
\end{enumerate}

\section{5. Aplicación Práctica: Evaluación de Perfiles de
Clientes}\label{aplicaciuxf3n-pruxe1ctica-evaluaciuxf3n-de-perfiles-de-clientes}

Utilizando la calculadora de churn, estime la probabilidad de abandono
para los siguientes casos:

\subsection{Caso 1: Cliente poco
vinculado}\label{caso-1-cliente-poco-vinculado}

\begin{itemize}
\tightlist
\item
  CreditScore: 650\\
\item
  Geography: France\\
\item
  Gender: Female\\
\item
  Age: 40\\
\item
  Tenure: 5\\
\item
  Balance: 100000\\
\item
  NumOfProducts: 1\\
\item
  HasCrCard: 0\\
\item
  IsActiveMember: 0\\
\item
  EstimatedSalary: 100000
\end{itemize}

Probabilidad estimada: \_\_\_\_\_\_\_\_\_\_\_\_\_\_\_\_\_\_\_

\subsection{Caso 2: Cliente joven y
activo}\label{caso-2-cliente-joven-y-activo}

\begin{itemize}
\tightlist
\item
  CreditScore: 720\\
\item
  Geography: Spain\\
\item
  Gender: Male\\
\item
  Age: 30\\
\item
  Tenure: 6\\
\item
  Balance: 50000\\
\item
  NumOfProducts: 2\\
\item
  HasCrCard: 1\\
\item
  IsActiveMember: 1\\
\item
  EstimatedSalary: 80000
\end{itemize}

Probabilidad estimada: \_\_\_\_\_\_\_\_\_\_\_\_\_\_\_\_\_\_\_

\subsection{Caso 3: Cliente mayor con alto saldo e
inactivo}\label{caso-3-cliente-mayor-con-alto-saldo-e-inactivo}

\begin{itemize}
\tightlist
\item
  CreditScore: 580\\
\item
  Geography: Germany\\
\item
  Gender: Female\\
\item
  Age: 55\\
\item
  Tenure: 3\\
\item
  Balance: 150000\\
\item
  NumOfProducts: 1\\
\item
  HasCrCard: 1\\
\item
  IsActiveMember: 0\\
\item
  EstimatedSalary: 90000
\end{itemize}

Probabilidad estimada: \_\_\_\_\_\_\_\_\_\_\_\_\_\_\_\_\_\_\_

\subsection{Caso 4: Cliente altamente
vinculado}\label{caso-4-cliente-altamente-vinculado}

\begin{itemize}
\tightlist
\item
  CreditScore: 700\\
\item
  Geography: France\\
\item
  Gender: Male\\
\item
  Age: 45\\
\item
  Tenure: 9\\
\item
  Balance: 120000\\
\item
  NumOfProducts: 3\\
\item
  HasCrCard: 1\\
\item
  IsActiveMember: 1\\
\item
  EstimatedSalary: 110000
\end{itemize}

Probabilidad estimada: \_\_\_\_\_\_\_\_\_\_\_\_\_\_\_\_\_\_\_

\section{6. Análisis Comparativo
Final}\label{anuxe1lisis-comparativo-final}

\begin{enumerate}
\def\labelenumi{\arabic{enumi}.}
\setcounter{enumi}{10}
\item
  ¿Cuál de los cuatro clientes presenta mayor riesgo de abandono?
\item
  ¿Qué variables parecen explicar las diferencias observadas?
\item
  Desde el punto de vista estratégico, ¿a cuál cliente priorizaría para
  una campaña de retención?
\end{enumerate}

\section{7. Matriz de Confusión y Métricas de
Clasificación}\label{matriz-de-confusiuxf3n-y-muxe9tricas-de-clasificaciuxf3n}

\begin{enumerate}
\def\labelenumi{\arabic{enumi}.}
\setcounter{enumi}{13}
\tightlist
\item
  A partir de la matriz de confusión mostrada en la aplicación,
  identifique los valores de:
\end{enumerate}

\begin{itemize}
\tightlist
\item
  TP (Verdaderos Positivos)\\
\item
  TN (Verdaderos Negativos)\\
\item
  FP (Falsos Positivos)\\
\item
  FN (Falsos Negativos)
\end{itemize}

Explique qué representa cada uno en el contexto de predicción de churn.

\begin{enumerate}
\def\labelenumi{\arabic{enumi}.}
\setcounter{enumi}{14}
\tightlist
\item
  Con base en los valores anteriores, calcule manualmente los siguientes
  indicadores utilizando sus fórmulas:
\end{enumerate}

Exactitud (Accuracy): Accuracy = (TP + TN) / (TP + TN + FP + FN)

Sensibilidad (Recall, clase Churn): Sensibilidad = TP / (TP + FN)

Especificidad: Especificidad = TN / (TN + FP)

Interprete cada resultado en el contexto del modelo de predicción de
abandono y compare con los valores reportados por la aplicación.

\section{Objetivo del Taller}\label{objetivo-del-taller}

Este ejercicio busca que el estudiante:

\begin{itemize}
\tightlist
\item
  Muestre e interprete distribuciones de variables relevantes.
\item
  Analice relaciones bivariadas con la variable respuesta.
\item
  Identifique factores de riesgo (OR \textgreater{} 1) y factores
  protectores (OR \textless{} 1) a partir de los Odds Ratios.
\item
  Relacione resultados estadísticos con decisiones de negocio.
\end{itemize}

\end{document}
